In this section, a brief overview of Generative Adversarial Network (GAN) is presented. The fundamental concept of the minmax game in GAN is introduced along with a discussion of different GAN architectures, such as the deep conditional network used in DCGAN or the multilayer-perceptrons employed in Vanilla GAN. Additionally, this section provides an introduction to Wasserstein GAN and Wasserstein GAN with Gradient Punishment, which use the Wasserstein distance as a loss function instead of the traditional binary cross-entropy loss.

\subsection{Core ideas}
\subsubsection{Basic GAN}
Generative Adversarial Network (GAN) are a type of deep learning algorithm that can be used to generate new data that mimics the characteristics of a given dataset. GANs consist of two neural networks, a generator network and a discriminator network, that are trained together in a two-player minimax game. Figure \ref{fig:gan_structure} illustrates an example of a GAN structure.  

During training, the generator and discriminator models work together. The generator generates a batch of samples, which are then combined with real examples from the domain and presented to the discriminator to be classified as either real or fake. The discriminator is updated to improve its ability to distinguish between real and fake samples in the next round. Importantly, the generator is updated based on how successful it was at fooling the discriminator with its generated samples. In this case, zero-sum means that when the discriminator successfully identifies real and fake samples, it is rewarded or no change is needed to the model parameters, whereas the generator is penalized with large updates to model parameters. Alternately, when the generator fools the discriminator, it is rewarded, or no change is needed to the model parameters, but the discriminator is penalized and its model parameters are updated. 

\begin{figure}[!h]
	 \centerline{\includegraphics[width=0.5\textwidth]{src/background/gan_structure.png}}
	 {\caption{Example of the Generative Adversarial Network Model Architecture \cite{ganstructure}}\label{fig:gan_structure}}
  
\end{figure}

\subsubsection{Conditional Generative Adversarial Networks (cGAN)}
A conditional Generative Adversarial Network  (cGAN) is a type of GAN that includes additional information in the input. Instead of just using noise as input for the generator, a conditional GAN appends a one-hot vector that indicates the class of the input. Similarly, for the discriminator, one-hot vectors are appended to the picture to predict the class as well. This allows the generator to generate samples conditioned on specific classes, resulting in more realistic and targeted samples. Figure \ref{fig:congan_structure} illustrates the example of a conditional GAN. 

\begin{figure}
	 \centerline{\includegraphics[width=0.8\textwidth]{src/background/congan_structure.png}}
	 {\caption{Example of the conditional Generative Adversarial Network Model Architecture \cite{ganstructure}}\label{fig:congan_structure}}
\end{figure}


\subsection{GAN Implementations}
The Vanilla GAN architecture, which solely consists of multi-layer perceptrons, often yields low-resolution images. In contrast, DCGAN addresses this issue by utilizing convolutional neural networks (CNNs) for both the generator and discriminator. Furthermore, the Vanilla GAN suffers from mode collapse, where the generator produces a limited range of samples that fail to represent the entire distribution of the training data. To overcome this limitation, Wasserstein distance is used as the loss function instead of the traditional cross-entropy loss (Kullback-Leibler divergence). In this section, we will provide an overview of these improvements.
\subsubsection{DCGAN}
\begin{figure}[!h]
	 \centerline{\includegraphics[width=0.8\textwidth]{src/background/dcgan.png}}
	 {\caption{Architecture of DCGAN Generator \cite{radford2015unsupervised}}\label{fig:dcgan}}
\end{figure}

DCGAN \cite{radford2015unsupervised} is a GAN architecture that uses convolutional neural networks (CNNs) as both the generator and discriminator networks to generate high-quality images. It follows a set of guidelines, such as using strided convolutions instead of pooling layers, batch normalization in both generator and discriminator, and ReLU activation in the generator with Tanh for the output, and LeakyReLU activation in the discriminator for all layers. Figure \ref{fig:dcgan} shows the architecture of DCGAN generator.


\subsubsection{WGAN}
Wasserstein GAN (WGAN) is a type of GAN that uses Wasserstein distance instead of traditional loss functions like cross-entropy. Wasserstein distance, also known as earth mover's distance, is a measure of distance between two probability distributions. It measures the amount of "work" or "distance" required to transform one distribution into the other. In other words, it quantifies the dissimilarity between two distributions. This not only mitigates the problem of mode collapse, where the generator produces a limited range of samples, but also establishes a reliable stopping point for training when the distance falls below a certain threshold.

Figure \ref{fig:wasserstein_distance} provides a demonstration of the vanishing gradient problem in regular GANs and the effectiveness of Wasserstein distance in addressing this issue \cite{arjovsky2017wasserstein}. It reveals that the gradient of the GAN discriminator becomes unstable and either explodes or diminishes, which negatively affects the learning process of the generator. In contrast, the WGAN critic is trained using Wasserstein distance as the loss function, leading to a smoother and more stable gradient.

\begin{figure}[!h]
	 \centerline{\includegraphics[width=0.8\textwidth]{src/background/wasserstein_distance.png}}
	 {\caption{Illustration of Wasserstein Distance\cite{arjovsky2017wasserstein}}\label{fig:wasserstein_distance}}
\end{figure}

\subsubsection{WGAN-GP}
To enforce the Lipschitz constraint (the gradients do not explode during training), weight clipping is used in WGAN. However, one of the major problems with weight clipping is that it can lead to vanishing gradients, which can slow down the training process and result in poor performance. Figure \ref{fig:weight_clipping} illustrates this hyperparamter-sensitive problem.

\begin{figure}[!h]
	 \centerline{\includegraphics[width=0.8\textwidth]{src/background/weight_clipping.png}}
	 {\caption{Illustration of the sensitivity problem of gradient to weight clipping in regular WGAN \cite{gulrajani2017improved}}\label{fig:weight_clipping}}
\end{figure}

Wasserstein GAN with Gradient Penalty (WGAN-GP) is a modification to the original WGAN proposed to address the issue of gradient explosions or vanishing gradients during the training of WGANs. The WGAN-GP method replaces the weight clipping constraint used in WGAN with a gradient penalty term that enforces the Lipschitz continuity of the critic function. Specifically, WGAN-GP penalizes the model if the gradient norm moves away from its target norm value 1.


